%===============================================================================
% LLMWiki Technical Specification — Volume 05: Agent Runtime
%===============================================================================
\chapter{Agent Runtime}
\label{cha:agent-runtime}

This volume specifies the \emph{agent runtime}---the execution environment that carries out query plans, orchestrates tool calls, manages conversation state, and streams results. The runtime is the bridge between the Query Planner (Vol.~07) and the external world (API, CLI, LLM providers).

\section{Design Goals}
\label{sec:runtime:goals}

\begin{designprinciple}[Stateless Execution]
Each query execution is a self-contained episode. No persistent agent state across queries; all context is passed explicitly in the plan and conversation history.
\end{designprinciple}

\begin{designprinciple}[Tool-First Architecture]
LLMs are not the primary actors; tools are. The LLM plans and decides; tools execute deterministically. This inverts the typical ``agent = LLM + tools'' model.
\end{designprinciple}

\begin{designprinciple}[Streaming by Default]
All long-running operations (search, synthesis, tool calls) stream partial results. The API responds with Server-Sent Events (SSE) or WebSocket frames.
\end{designprinciple}

\begin{designprinciple}[Budget Enforcement]
Every execution has configurable budgets: token budget, time budget, tool call budget, cost budget. Exceeding any budget terminates with a partial result.
\end{designprinciple}

\begin{designprinciple}[Observability Native]
Every decision (plan step, tool choice, retrieval result, synthesis token) is logged with structured traces. OpenTelemetry compatible.
\end{designprinciple}

\section{Runtime Architecture}
\label{sec:runtime:architecture}

\begin{figure}[htbp]
\centering
\begin{tikzpicture}[
  mod/.style={draw, rounded corners, fill=llmwikiorange!10, minimum width=3cm, minimum height=1.2cm, font=\small},
  flow/.style={->, thick, >=stealth}
]
  \node[mod] (planner) at (0,2) {Query Planner};
  \node[mod] (executor) at (4,2) {Plan Executor};
  \node[mod] (toolreg) at (8,2) {Tool Registry};
  \node[mod] (llm) at (8,0) {LLM Client};
  \node[mod] (state) at (4,0) {State Manager};
  \node[mod] (stream) at (0,0) {Stream Handler};

  \draw[flow] (planner) -- (executor);
  \draw[flow] (executor) -- (toolreg);
  \draw[flow] (executor) -- (state);
  \draw[flow] (executor) -- (llm);
  \draw[flow] (executor) -- (stream);
  \draw[flow] (state) -- (stream);
  \draw[flow] (llm) -- (stream);
  \draw[flow] (toolreg) -- (stream);
\end{tikzpicture}
\caption{Agent Runtime components.}
\label{fig:runtime:arch}
\end{figure}

\subsection{Core Components}
\label{sec:runtime:components}

\begin{itemize}
  \item \textbf{Plan Executor}: Traverses the query plan DAG, executes nodes in topological order, handles parallelism and dependencies.
  \item \textbf{Tool Registry}: Maps tool names to implementations, handles invocation, timeouts, retries, and result normalization.
  \item \textbf{LLM Client}: Abstracts LLM providers (local: llama.cpp, vLLM, ollama; remote: OpenAI, Anthropic, etc.) with unified streaming interface.
  \item \textbf{State Manager}: Maintains scratchpad, conversation history, intermediate results, and budget counters.
  \item \textbf{Stream Handler}: Emits SSE/WebSocket events for partial results, progress, errors, and final answer.
\end{itemize}

\section{Plan Execution Model}
\label{sec:runtime:plan-exec}

\subsection{Query Plan Structure}
\label{sec:runtime:plan-structure}

The planner (Vol.~07) produces a \texttt{QueryPlan}:

\begin{verbatim}
struct QueryPlan {
  id: PlanId,
  query: String,
  steps: Vec<PlanStep>,
  edges: Vec<(StepId, StepId)>,  // DAG dependencies
  budget: ExecutionBudget,
  config: PlanConfig,
}

struct PlanStep {
  id: StepId,
  kind: StepKind,        // Search, Synthesize, ToolCall, Aggregate, Branch
  tool: Option<ToolName>,
  input: StepInput,      // Template with variable interpolation
  output: StepOutput,    // Expected output schema
  depends_on: Vec<StepId>,
  parallel: bool,
  retry: RetryPolicy,
  timeout_ms: u64,
}
\end{verbatim}

\subsection{Execution Algorithm}
\label{sec:runtime:exec-algo}

\begin{algorithm}[H]
\caption{Plan Execution}
\label{alg:plan-exec}
\begin{algorithmic}[1]
\Require Plan $P$, initial state $S_0$
\Ensure Final state $S_f$ with answer
\State $S \gets S_0$, $completed \gets \varnothing$
\State $ready \gets \{ s \in P.steps \mid s.depends\_on \subseteq completed \}$
\While{$ready \neq \varnothing$}
  \If{$\exists s \in ready: s.parallel$}
    \State $batch \gets \{ s \in ready \mid s.parallel \}$
    \State $results \gets \textsc{ExecuteParallel}(batch, S)$
  \Else
    \State $s \gets \textsc{Pop}(ready)$
    \State $results \gets \textsc{ExecuteStep}(s, S)$
  \EndIf
  \State $S \gets \textsc{UpdateState}(S, results)$
  \State $completed \gets completed \cup \{s.id \mid s \in batch \lor s\}$
  \State $ready \gets \{ s \in P.steps \setminus completed \mid s.depends\_on \subseteq completed \}$
  \State \textsc{CheckBudget}(S, P.budget) \Comment{May raise BudgetExceeded}
  \State \textsc{EmitProgress}(S, completed) \Comment{Stream update}
\EndWhile
\State \Return $S$
\end{algorithmic}
\end{algorithm}

\subsection{Step Kinds}
\label{sec:runtime:step-kinds}

\begin{itemize}
  \item \texttt{Search}: Invoke hybrid search + evidence gate (Vol.~04)
  \item \texttt{Synthesize}: Call LLM with evidence context, stream tokens
  \item \texttt{ToolCall}: Invoke registered tool (see \S\ref{sec:runtime:tools})
  \item \texttt{Aggregate}: Merge results from parallel branches
  \item \texttt{Branch}: Conditional execution based on prior result
\end{itemize}

\section{Tool System}
\label{sec:runtime:tools}

\subsection{Tool Trait}
\label{sec:runtime:tool-trait}

\begin{verbatim}
trait Tool {
    fn name(&self) -> &str;
    fn description(&self) -> &str;
    fn schema(&self) -> ToolSchema;          // JSON Schema for input
    fn output_schema(&self) -> ToolSchema;   // JSON Schema for output
    fn execute(&self, input: ToolInput) -> Result<ToolOutput>;
    fn execute_stream(&self, input: ToolInput) -> Stream<ToolOutput>;
    fn timeout_ms(&self) -> u64 { 30000 }
    fn retry(&self) -> RetryPolicy { RetryPolicy::default() }
}
\end{verbatim}

\subsection{Built-in Tools}
\label{sec:runtime:builtin-tools}

\begin{table}[htbp]
\centering
\caption{Built-in Tools}
\label{tab:runtime:builtin-tools}
\begin{tabular}{lll}
\toprule
\textbf{Tool} & \textbf{Input} & \textbf{Output} \\
\midrule
\texttt{search} & \texttt{SearchInput}\{query, k, filters\} & \texttt{SearchOutput}\{results[], citations[]\} \\
\texttt{graph\_traverse} & \texttt{TraverseInput}\{start, pattern, hops\} & \texttt{Subgraph}\{nodes[], edges[]\} \\
\texttt{read\_unit} & \texttt{UnitId} & \texttt{SemanticUnit} \\
\texttt{read\_file} & \texttt{FilePath} & \texttt{FileContent} \\
\texttt{llm\_complete} & \texttt{CompletionInput}\{prompt, params\} & \texttt{Stream<String>} \\
\texttt{extract\_entities} & \texttt{ExtractInput}\{text, types\} & \texttt{Entities}\{entities[], relations[]\} \\
\texttt{summarize} & \texttt{SummarizeInput}\{units[], max\_len\} & \texttt{String} \\
\texttt{verify\_evidence} & \texttt{VerifyInput}\{query, units[]\} & \texttt{VerifiedUnits}\{verified[]\} \\
\bottomrule
\end{tabular}
\end{table}

\subsection{Custom Tools}
\label{sec:runtime:custom-tools}

Users register custom tools via config:

\begin{verbatim}
[tools.custom.my_tool]
module = "my_tools"
class = "MyTool"
timeout_ms = 10000
retry = { max_attempts = 2, backoff_ms = 500 }
\end{verbatim}

Tools can be synchronous (return \texttt{Result}) or streaming (return \texttt{Stream}). Streaming tools emit partial results for long operations (e.g., file scanning, API pagination).

\section{LLM Client Abstraction}
\label{sec:runtime:llm-client}

\begin{verbatim}
trait LLMClient {
    fn complete(&self, req: CompletionRequest) -> Result<CompletionResponse>;
    fn stream(&self, req: CompletionRequest) -> Stream<CompletionChunk>;
    fn embed(&self, texts: Vec<String>) -> Result<Vec<Vec<f32>>>;
    fn models(&self) -> Vec<ModelInfo>;
    fn count_tokens(&self, text: &str) -> usize;
}
\end{verbatim}

\paragraph{Providers:}
\begin{itemize}
  \item \textbf{Local}: llama.cpp (server), vLLM, ollama, TGI
  \item \textbf{Remote}: OpenAI, Anthropic, Azure, Vertex, Bedrock
  \item \textbf{Unified API}: OpenAI-compatible endpoint
\end{itemize}

The runtime selects the provider per plan step (e.g., small local model for planning, large model for synthesis).

\section{State Management}
\label{sec:runtime:state}

\begin{verbatim}
struct ExecutionState {
  plan_id: PlanId,
  step_results: HashMap<StepId, StepResult>,
  scratchpad: Scratchpad,           // Key-value store for LLM use
  conversation: Vec<Message>,       // User/assistant messages
  evidence: Vec<VerifiedUnit>,      // Accumulated evidence
  budget: BudgetTracker,            // Tokens, time, cost, calls
  trace: TraceContext,              // OpenTelemetry trace
}
\end{verbatim}

\paragraph{Scratchpad:} Structured key-value store accessible to LLM via tool calls. Used for intermediate reasoning, variable passing between steps, and user-facing ``thoughts''.

\paragraph{Budget Tracker:} Hard limits with configurable actions on exceed:
\begin{verbatim}
budget:
  max_tokens: 100000
  max_time_ms: 120000
  max_tool_calls: 50
  max_cost_usd: 0.50
  on_exceed: "truncate"  # "error" | "truncate" | "partial"
\end{verbatim}

\section{Streaming Protocol}
\label{sec:runtime:streaming}

API uses Server-Sent Events (SSE) over HTTP or WebSocket:

\begin{verbatim}
# Event types
event: plan_start     { "plan_id": "...", "steps": N }
event: step_start     { "step_id": "...", "kind": "Search" }
event: step_progress  { "step_id": "...", "progress": 0.5, "details": "..." }
event: evidence       { "unit_id": "...", "relevance": 0.94, "preview": "..." }
event: token          { "token": "the", "step_id": "..." }
event: step_complete  { "step_id": "...", "output": {...} }
event: plan_complete  { "answer": "...", "citations": [...], "usage": {...} }
event: error          { "step_id": "...", "error": "..." }
event: budget_warning { "type": "tokens", "used": 80000, "limit": 100000 }
\end{verbatim}

Clients can render progressive UI: show plan, show evidence as it arrives, stream answer tokens.

\section{Error Handling and Retries}
\label{sec:runtime:errors}

\begin{itemize}
  \item \textbf{Transient errors} (network, timeout): Retry per step policy (exponential backoff)
  \item \textbf{Tool errors}: Captured in \texttt{StepResult.error}, execution continues if step not critical
  \item \textbf{LLM errors}: Rate limit $\to$ backoff; context overflow $\to$ truncate/summarize history
  \item \textbf{Budget exceeded}: Emit \texttt{budget\_warning}, then \texttt{plan\_complete} with partial answer
  \item \textbf{Fatal errors}: Invalid plan, OOM, provider unavailable $\to$ abort with error event
\end{itemize}

\section{Conversation Management}
\label{sec:runtime:conversation}

Multi-turn conversations maintain history in \texttt{ExecutionState.conversation}. The planner can reference prior turns for context. History is pruned to fit budget (summarize old turns, keep recent).

\section{Patent-Relevant Novelty}
\label{sec:runtime:patent}

\begin{patentclaim}
\textbf{Claim 8:} An agent runtime wherein: (a) query plans are DAGs executed by a deterministic executor with topological scheduling, (b) tools are first-class citizens with JSON Schema contracts, streamable execution, and configurable timeouts/retries, (c) LLM invocation is abstracted behind a provider-agnostic client supporting both local and remote models, (d) execution state includes a structured scratchpad accessible to tools and LLM, and (e) hard budgets (tokens, time, cost, calls) are enforced at the executor level with configurable overflow behavior.
\end{patentclaim}

\section{Summary}
\label{sec:runtime:summary}

The Agent Runtime executes query plans as deterministic, budgeted, observable tool orchestrations:
\begin{enumerate}
  \item \textbf{Plan executor}: DAG-aware, parallel-capable, topological
  \item \textbf{Tool system}: Schema-contracted, streamable, extensible
  \item \textbf{LLM abstraction}: Multi-provider, local-first, token-aware
  \item \textbf{State}: Scratchpad, evidence accumulator, budget tracker
  \item \textbf{Streaming}: SSE/WebSocket with fine-grained event types
  \item \textbf{Budgets}: Hard limits with graceful degradation
\end{enumerate}
The runtime is the ``CPU'' of LLMWiki; the Query Planner (Vol.~07) is the compiler that generates its instruction sequences.

%===============================================================================
% End of Volume 05
%===============================================================================
